\chapter{Conclusions}
\label{ch:conclusions}

This thesis has investigated how computational approaches can enhance both the strategic design and the daily operation of docked \acp{bss}. By combining spatial optimization, mobility and battery modeling, and predictive maintenance, the research addresses complementary challenges across the life cycle of \acp{bss}. Applied to Barcelona and its \ac{bss}, \textit{Bicing}, the work demonstrates how operational and contextual data can be transformed into actionable insights that support more efficient, equitable, and sustainable shared mobility systems.

The first major contribution of this thesis is a flexible and modular framework for station planning (Chapter~\ref{ch:station_planning_in_barcelona}). By integrating \ac{mcdm} scoring, directed network representations, and topography-adjusted distances, the framework enables candidate station locations to be evaluated with greater spatial realism and granularity than in previously existing approaches. It supports the explicit exploration of trade-offs among competing priorities, including demand, first--last-kilometer connectivity, social equity, and infrastructure integration, while accounting for inter-station proximity, system accessibility, and station density constraints. The use of a \ac{ga} demonstrates the suitability of metaheuristic methods for navigating large urban search spaces and producing robust, context-sensitive station configurations that support incremental network expansion.

The second contribution introduces a joint prediction framework for \ac{e-b} mobility and battery levels, combining a Markov-chain mobility model with a physics-informed battery dynamics approach (Chapter~\ref{ch:results_mobility_and_battery_modeling}). The framework is grounded in the mobility analysis presented in Chapter~\ref{ch:bicing_mobility_patterns} and the battery analysis in Section~\ref{subsec:battery_results_analysis}, which together identify the main drivers of mode-specific mobility and energy depletion. These insights are incorporated into the model through determinants such as trip distance, elevation, and temporal regularities. In application, the framework supports system-level forecasts of charging needs and potential shortages, enabling more proactive rebalancing and charging strategies.

The third contribution focuses on \ac{pm} modeling using \ac{sa} and \ac{ml} techniques (Chapter~\ref{ch:results_failure_forecasting}). It builds on the mobility analysis in Chapter~\ref{ch:bicing_mobility_patterns} and the maintenance data analysis in Section~\ref{subsec:MOs_analysis}. By combining historical mobility indicators, weather-related variables, and detailed maintenance records, the thesis develops interpretable models that estimate component-level failure risks for both \ac{m-b} and \ac{e-b}. The results show that degradation is closely associated with usage intensity and bike model type. In practice, the models support proactive workshop scheduling, improved spare-part planning, and reduced downtime by anticipating failures before they occur.

Taken together, these contributions address station planning, operational forecasting, and maintenance as distinct but related components of a \ac{bss}. While each component is modeled separately, the thesis highlights conceptual links across these stages through shared data, mobility patterns, and structural characteristics of the city. Moreover, although Barcelona and \textit{Bicing} served as the empirical case studies, the methods developed in this thesis are broadly transferable to other urban contexts. The station planning framework can be adapted to different urban forms and policy priorities, while the mobility, battery, and maintenance models rely on data sources commonly available to \ac{bss} operators.

In conclusion, this thesis demonstrates that integrating computational approaches into the planning and operation of \acp{bss} is both feasible and beneficial. By advancing methodological frameworks and empirically illustrating their potential using data from one of Europe's largest \acp{bss}, this research has contributed to the development of more efficient, resilient, and user-centered shared mobility systems and supports the broader transition toward sustainable and intelligent urban mobility.